\section{Referencia rápida de términos y métricas}

\subsection{Tabla de correspondencia de términos clave}
\begin{longtable}{p{2.7cm}p{3.2cm}p{5.6cm}}
\toprule
\textbf{Término} & \textbf{Comprensión en español} & \textbf{Función en esta clase} \\
\midrule
\endfirsthead
\toprule
\textbf{Término} & \textbf{Comprensión en español} & \textbf{Función en esta clase} \\
\midrule
\endhead
API Access & Acceso de caja negra a la interfaz & Permite experimentación a nivel de comportamiento, pero dificulta la atribución a nivel de mecanismo. \\
Open Weights & Pesos abiertos & Permite el análisis interno, el fine-tuning y la destilación, pero sigue restringido por la receta de entrenamiento existente. \\
Open Source & Apertura de toda la cadena & Abarca datos, modelo, entrenamiento y evaluación, y permite la reproducibilidad a nivel de sistema. \\
Access Shapes Research & El acceso determina la forma de la investigación & Hilo conductor del curso: los recursos accesibles determinan qué tipo de preguntas se pueden plantear. \\
Agent Controller & Controlador de ejecución & El LLM no es solo un módulo al que se invoca, sino el núcleo de control del flujo. \\
Memory Stream & Flujo de memoria & Registra observaciones y acciones históricas, y da contexto para la recuperación y la reflexión posteriores. \\
Retrieval & Recuperación & Extrae del historial de memoria el estado más relevante para la tarea actual. \\
Reflection & Reflexión & Comprime la trayectoria histórica en una estrategia de alto nivel, lo que reduce la repetición de errores. \\
Planning & Planificación & Descompone una tarea larga en subobjetivos verificables. \\
Tool Use & Uso de herramientas & Convierte las decisiones del lenguaje en acciones de ejecución reales. \\
Orchestration & Orquestación & Organiza los límites de roles, los límites de contexto y el orden de ejecución de múltiples agentes. \\
Benchmark & Evaluación de referencia & Convierte la impresión de «parece más fuerte» en un resultado comparable y verificable. \\
Failure Taxonomy & Taxonomía de fallas & Distingue tipos de falla como malentendido del objetivo, falla de herramienta o degradación de la estrategia. \\
Replayability & Capacidad de repetición & Permite que cada falla se pueda repetir y diagnosticar; es la base de la reproducibilidad. \\
Audit Log & Registro de auditoría & Aporta la cadena de evidencia para la seguridad y la rendición de cuentas. \\
Distribution Shift & Cambio de distribución & Evalúa la estabilidad y la capacidad de generalización del agente en escenarios fuera de distribución. \\
Safety Gate & Compuerta de seguridad & Establece confirmación de permisos y mecanismos de reversión para acciones de alto riesgo. \\
Reproducibility Package & Paquete de reproducibilidad & Reúne configuración, scripts, entorno y resultados para que un tercero los verifique. \\
Open Definition & Definición de apertura & Precisa los límites de «apertura» para evitar que el uso confuso del término distorsione la colaboración. \\
Research Throughput & Rendimiento de la investigación & Lo determinan conjuntamente el cómputo, la cadena de herramientas y los procesos organizacionales. \\
\bottomrule
\end{longtable}

\begin{knowledgebox}{Por qué hacer una referencia rápida de términos}
La investigación de agentes abarca modelos, sistemas, evaluación y gobernanza, y el vocabulario que usan distintos equipos no siempre coincide.
Estandarizar los términos puede reducir de manera significativa los malentendidos en la colaboración, sobre todo cuando se reproducen experimentos entre instituciones.
\end{knowledgebox}

\subsection{Resumen de la sección}
Unificar la terminología es una acción de ingeniería de bajo costo y alto rendimiento. No mejora directamente el puntaje en las tablas de clasificación, pero sí mejora de manera significativa la comparabilidad de los experimentos del equipo y la velocidad de acumulación de conocimiento.

